%\documentclass[12pt,oneside]{amsart}
\documentclass{scrartcl}

%\documentclass{article}

\usepackage{amsthm,amsmath,amssymb}
\usepackage[numbers]{natbib}

%\usepackage[colorlinks]{hyperref}
\usepackage{hyperref}
\hypersetup{
    colorlinks,%
    citecolor=black,%
    filecolor=black,%
    linkcolor=black,%
    urlcolor=black
}
\usepackage{hypernat}
\usepackage[top=2.5cm, bottom=2.5cm, left=2.8cm,
right=2.8cm]{geometry}
\usepackage{graphicx}


%\setlength{\parskip}{0pt}
%\setlength{\parsep}{0pt}
%\setlength{\headsep}{0pt}
%\setlength{\topskip}{0pt}
%\setlength{\topmargin}{0pt}
%\setlength{\topsep}{0pt}
%\setlength{\partopsep}{0pt}

\linespread{1}

%\usepackage{marvosym}

%\textwidth = 6.5 in \textheight = 9.2 in \oddsidemargin = 0.0 in
%\evensidemargin = 0.0 in \topmargin = -0.0 in \headheight = 0.0 in
%\headsep = 0.0in \parskip = 0.0in \parindent = 0.0in
%\def\baselinestretch{0.871}


\newtheorem{theorem}{Theorem}[section]
\newtheorem{proposition}[theorem]{Proposition}
\newtheorem{problem}[theorem]{Problem}
\newtheorem{fact}[theorem]{Fact}
\newtheorem{lemma}[theorem]{Lemma}
\newtheorem{defn}[theorem]{Definition}
\newtheorem{corollary}[theorem]{Corollary}
\newtheorem{remark}{Remark}[section]
\newtheorem{example}[theorem]{Example}

\newcommand{\E}{{\mathbb E}}
\newcommand{\se}{{\mathcal{E}}}
\newcommand{\R}{{\mathbb R}}
\renewcommand{\P}{{\mathbb P}}
\newcommand{\ac}{{\mathcal{A}}}
\newcommand{\A}{{\mathcal{A}}}
\newcommand{\B}{{\mathcal{B}}}
\newcommand{\C}{{\mathcal{C}}}
\newcommand{\M}{{\mathcal{M}}}
\newcommand{\Mf}{{\mathfrak{M}}}
\newcommand{\T}{{\mathcal{T}}}
\newcommand{\Z}{{\mathcal{Z}}}
\newcommand{\K}{{\mathcal{K}}}
\newcommand{\lc}{{\mathcal{L}}}
\newcommand{\Ps}{{\mathcal{P}}}
\newcommand{\G}{{\mathcal{G}}}
\newcommand{\pa}{{\mathring{p}}}
\newcommand{\F}{{\cal F}}
\newcommand{\samp}{{\mathcal{X}}}
\newcommand{\X}{{\mathcal{X}}}
\newcommand{\hi}{{\hat{\Phi}}}
\newcommand{\data}{{\text{data}}}

\newcounter{rcnt}[section]
\renewcommand{\thercnt}{(\roman{rcnt})}

\def\qt#1{\qquad\text{#1}}

\def\argmin{\mathop{\rm argmin}}
\def\argmax{\mathop{\rm argmax}}


\setlength{\parskip}{1.4 \medskipamount}
%\sloppy
%\linespread{1.3}

\begin{document}

\title{STAT 238 - Bayesian Statistics  \\
  Lecture Sixteen} 
\subtitle{Spring 2026, UC Berkeley}

\author{Aditya Guntuboyina}


\date{02 Feb 2026}

\maketitle

\section{Linear Regression}

In the last lecture, we discussed Bayesian inference for multiple
linear regression. The setting is as follows: we have one
response variable $y$ and 
$m$ covariates $x_1, \dots, x_m$  ($m = 1$ corresponds to simple linear
regression). We observe data on $n$ instances or subjects for all
these variables: $(y_i, x_{i1}, \dots, x_{im})$ for $i = 1, \dots,
n$. The multiple linear regression model (with normal errors) is given
by: 
\begin{equation}\label{lm}
  y_i = \beta_0 + \beta_1 x_{i1} + \dots + \beta_m x_{im} + \epsilon_i
~~  \text{with $\epsilon_i \overset{\text{i.i.d}}{\sim} N(0, \sigma^2)$}. 
\end{equation}

Using the matrix vector notation: 
\begin{equation*}
  y =
  \begin{pmatrix}
    y_1 \\ \cdot \\ \cdot \\ \cdot \\ y_n
  \end{pmatrix} ~~~ X =
  \begin{pmatrix}
    1 & x_{11} & \dots & x_{1m} \\
    1 & x_{21} & \dots & x_{2m} \\
    \cdot & \cdot & \dots & \cdot \\
    \cdot & \cdot & \dots & \cdot \\
    \cdot & \cdot & \dots & \cdot \\
    1 & x_{n1} & \dots & x_{nm}
  \end{pmatrix} ~~~
  \beta =
  \begin{pmatrix}
    \beta_0 \\ \beta_1 \\ \cdot \\ \cdot \\ \cdot \\ \beta_m
  \end{pmatrix} ~~~
  \hat{\beta} =
  \begin{pmatrix}
    \hat{\beta}_0 \\ \hat{\beta}_1 \\ \cdot \\ \cdot \\ \cdot \\ \hat{\beta}_m
  \end{pmatrix}
\end{equation*}
the model can be written as:
\begin{align*}
  y = X \beta + \epsilon ~~ \text{ where } \epsilon \sim N(0, \sigma^2
  I_n). 
\end{align*}
Below we review standard frequentist and Bayesian approaches to linear
regression. Both lead to identical answers. 

\subsection{Frequentist Approach}

Here the focus is on the least squares estimate which is given by:
\begin{align*}
  \hat{\beta} = (X^T X)^{-1} X^T y. 
\end{align*}
To perform inference, the key is to compute the distribution of
$\hat{\beta}$. The dependence of the above on the $X$ matrix is quite
complicated. On the other hand, the dependence on $y$ is linear. If we
assume that $X$ is fixed (non-random), then it is easy to write the
distribution of $\hat{\beta}$. Indeed because $y \sim N(X \beta,
\sigma^2 I_n)$, it is easy to check that
\begin{align*}
  \hat{\beta} \sim N \left( \beta, \sigma^2 (X^T X)^{-1} \right). 
\end{align*}
This implies that
\begin{align*}
  \hat{\beta}_j \sim N(\beta_j, \sigma^2 (X^T X)^{j+1, j+1}) ~~ \text{
  for each } j = 0, 1, \dots, m, 
\end{align*}
where $(X^T X)^{j+1, j+1}$ is the $(j+1, j+1)$-th diagonal entry of
$(X^T X)^{-1}$. The above is equivalent to:
\begin{align}\label{sigk}
  \frac{\hat{\beta}_j - \beta_j}{\sigma \sqrt{(X^T X)^{j+1, j+1}}}
  \sim N(0, 1). 
\end{align}
Because $\sigma$ is unknown, it is natural to replace it with a
suitable estimate $\sigma$. The standard estimate used is the residual
standard error: 
\begin{equation*}
  \hat{\sigma} := \sqrt{\frac{S(\hat{\beta})}{n-m-1}}. 
\end{equation*}
where $S(\beta) = \|y - X \beta\|^2$ is the sum of squares and
$S(\hat{\beta})$ is the smallest possible value of the sum of
squares. Replacing $\sigma$ by $\hat{\sigma}$ in \eqref{sigk} will
change the distribution. In fact, it can be shown that $N(0, 1)$ will
be replaced by the standard $t$ distribution with $n-m-1$ degrees of
freedom:
\begin{align}\label{siguk}
  \frac{\hat{\beta}_j - \beta_j}{\hat{\sigma} \sqrt{(X^T X)^{j+1, j+1}}}
  \sim t_{n-m-1}. 
\end{align}
This can be used to derive confidence intervals for $\beta_j$. If 
  $t_{n-m-1, \alpha/2}$ is the point beyond which the $t$-distribution
  (with $n-m-1$ degrees of freedom) assigns probability $\alpha/2$,
  then
  \begin{equation*}
    \P \left\{-t_{n-m-1, \alpha/2} \leq     \frac{\beta_j - \hat{\beta}_j}{\hat{\sigma} \sqrt{(X^T X)^{j+1,
          j+1}}} \leq t_{n-m-1, \alpha/2}  \right\} = 
  1 - \alpha
\end{equation*}
which is same as:
\begin{equation*}
  \P \left\{\hat{\beta}_j - \hat{\sigma} \sqrt{(X^T X)^{j+1,
          j+1}} t_{n-m-1, \alpha/2} \leq \beta_j \leq \hat{\beta}_j - \hat{\sigma} \sqrt{(X^T X)^{j+1,
          j+1}} t_{n-m-1, \alpha/2}    \right\} = 1 -
    \alpha 
\end{equation*}
This interval:
\begin{equation}\label{freqint}
  \left[\hat{\beta}_j - \hat{\sigma} \sqrt{(X^T X)^{j+1,
          j+1}} t_{n-m-1, \alpha/2},  \hat{\beta}_j - \hat{\sigma} \sqrt{(X^T X)^{j+1,
          j+1}} t_{n-m-1, \alpha/2} \right]
  \end{equation}
  is the $100(1-\alpha)\%$ confidence interval for
  $\beta_j$.

The quantity $\hat{\sigma}
\sqrt{(X^T X)^{j+1, j+1}}$ is known as the standard error
corresponding to $\beta_j$.   

\subsection{Bayesian Approach}

The default prior for linear regression is: 
\begin{equation*}
  \beta_0, \beta_1, \dots, \beta_m, \log \sigma
  \overset{\text{i.i.d}}{\sim} \text{unif}(-\infty, \infty). 
\end{equation*}
The likelihood is given by:
\begin{align}\label{liklin}
  \propto \sigma^{-n} \exp \left(-\frac{1}{2 \sigma^2} \|y - X
  \beta\|^2 \right). 
\end{align}
To obtain this likelihood, we can assume that $X$ is fixed (and that
$\epsilon \sim N(0, \sigma^2 I_d)$). Or we can assume that $\epsilon$
is independent of $X$ and that the distribution of $X$ does not depend
on the parameters $\beta, \sigma$. There could be other assumptions on
$\epsilon, X$ which can lead to the same (or very similar) likelihood
(see Section \ref{armods}). 

Under this prior and likelihood, we calculated (see last lecture notes) that the
posterior density of $\beta$ is given by:
\begin{equation}\label{post3}
  \beta_0, \dots, \beta_m \mid \text{data} \sim
  t_{m+1}\left(\hat{\beta},\frac{S(\hat{\beta})}{n - m - 1} \left(X^T
      X 
  \right)^{-1}, n - m - 1 \right). 
\end{equation}
Note that, in \eqref{post3}, the quantity
$S(\hat{\beta})/(n-m-1)$ is the \textbf{frequentist 
  unbiased} estimator for $\sigma^2$, which is denoted by 
$\hat{\sigma}^2$. $\hat{\sigma}$ can also be justified as a Bayesian
estimator of $\sigma$ (this will be a question in Homework three). 

With the notation for $\hat{\sigma}$, the posterior \eqref{post3}
becomes:
\begin{equation}\label{post4}
  \beta_0, \dots, \beta_m \mid \text{data} \sim
  t_{m+1}\left(\hat{\beta},\hat{\sigma}^2 \left(X^T
      X 
  \right)^{-1}, n - m - 1 \right). 
\end{equation}
By one of the facts mentioned about the multivariate $t$-distribution,
the posterior of each individual $\beta_j$ is also $t$:  
  \begin{equation}\label{marg4}
    \beta_j \mid \text{data} \sim t_{1}\left(\hat{\beta}_j,
\hat{\sigma}^2 (X^T X)^{j+1, j+1}, n-m-1 \right) 
  \end{equation}
  where $(X^T X)^{j+1, j+1}$ is the $(j+1)^{th}$ diagonal entry 
  of $(X^T X)^{-1}$ (note that we are using the $(j+1)$th diagonal
  entry of $X^TX$ because $\beta_j$ is the $(j+1)$th component of
  $\beta$). This implies that
  \begin{equation}\label{bayes_t}
    \frac{\beta_j - \hat{\beta}_j}{\hat{\sigma} \sqrt{(X^T X)^{j+1,
          j+1}}} \mid \text{data} \sim \text{univariate standard } t \text{ with }
    n-m-1 \text{ d.f.}
  \end{equation}
  This can be used to obtain uncertainty intervals for $\beta_j$. If
  $t_{n-m-1, \alpha/2}$ is the point beyond which the $t$-distribution
  (with $n-m-1$ degrees of freedom) assigns probability $\alpha/2$,
  then
  \begin{equation*}
    \P \left\{-t_{n-m-1, \alpha/2} \leq     \frac{\beta_j - \hat{\beta}_j}{\hat{\sigma} \sqrt{(X^T X)^{j+1,
          j+1}}} \leq t_{n-m-1, \alpha/2} \mid \text{data} \right\} =
  1 - \alpha
\end{equation*}
which is same as:
\begin{equation*}
  \P \left\{\hat{\beta}_j - \hat{\sigma} \sqrt{(X^T X)^{j+1,
          j+1}} t_{n-m-1, \alpha/2} \leq \beta_j \leq \hat{\beta}_j - \hat{\sigma} \sqrt{(X^T X)^{j+1,
          j+1}} t_{n-m-1, \alpha/2}  \mid \text{data}  \right\} = 1 -
    \alpha 
\end{equation*}
This interval:
\begin{equation}\label{bayesint}
  \left[\hat{\beta}_j - \hat{\sigma} \sqrt{(X^T X)^{j+1,
          j+1}} t_{n-m-1, \alpha/2},  \hat{\beta}_j - \hat{\sigma} \sqrt{(X^T X)^{j+1,
          j+1}} t_{n-m-1, \alpha/2} \right]
  \end{equation}
  is the $100(1-\alpha)\%$ Bayesian Credible interval for
  $\beta_j$.

From \eqref{freqint} and \eqref{bayesint}, it is clear that the
Bayesian and frequentist $100(1
- \alpha)\%$ intervals for $\beta_j$ coincide.

\subsection{When $n$ is large}

The degrees of freedom corresponding to the $t$-density in \eqref{post3} is $n - m - 1$ where
$n$ is the number of observations, and $m$ is the number of
covariates. Thus \textbf{if $n-m-1$ is large}, then the posterior
distribution (which is actually $t$) is approximately normal: 
\begin{equation*}
    t_{m+1}\left(\hat{\beta},\frac{S(\hat{\beta})}{n - m - 1} \left(X^T
      X 
  \right)^{-1}, n - m - 1 \right) \approx N_{m+1}(\hat{\beta},
\frac{S(\hat{\beta})}{n - m - 1} \left(X^T X  \right)^{-1}).  
\end{equation*}
In other words, when $n-m-1$ is large, the $t$-density \eqref{post3} is approximately
equal to the $N_{m+1}(\hat{\beta}, \hat{\sigma}^2 (X^T
X)^{-1})$. Further, when $n-m-1$ is large, the distribution \eqref{marg4}
will be close to the normal  distribution $N(\hat{\beta}_j,
\hat{\sigma}^2 (X^T X)^{j+1, j+1})$. In such cases, $t_{n-m-1,
  \alpha/2}$ can be replaced by $z_{\alpha/2}$ in the intervals. 

\subsection{Bayesian vs Frequentist}
We mentioned above that the Bayesian uncertainty interval for
$\beta_j$ exactly coincides with the standard frequentist confidence
interval. But the Bayesian reasoning that led to the interval differs
fundamentally from the frequentist reasoning. The Bayesian reasoning
is based on the fact \eqref{bayes_t} while frequentist reasoning is
based on \eqref{siguk}. 

The only difference between \eqref{siguk} and \eqref{bayes_t} is in
the conditioning. In \eqref{bayes_t}, the random variable is $\beta_j$
and we are computing its distribution conditional on the observed
data. In contrast, in the frequentist statement \eqref{siguk}, the
random variable is $\hat{\beta}_j$ (which is a function of $y_1,
\dots, y_n$) and we are computing its distribution the parameters
$\beta, \sigma$ are fixed.  Note that $X$ is assumed to be fixed
here.

The fact that both Bayesian and frequentist reasoning lead to the same
interval is a matter of coincidence. By changing the setting slightly,
this coincidence can be broken. One such example involves AR models
and is discussed next. 

\section{AutoRegression}\label{armods}

We observe a time series $\{y_t, t = 1, \dots, n\}$. The
AutoRegressive Model of order 1 (in short, AR(1)) is given by:
\begin{align}\label{ar1}
  y_t = \beta_0 + \beta_1 y_{t-1} + \epsilon_t ~~\text{ where $t = 2,
  \dots, n$}. 
\end{align}
In the above equation, $y_1$ does not appear on the right hand side,
so we can treat $y_1$ as a constant. The equation \eqref{ar1} is
simply a linear regression model with covariate given by $x_t =
y_{t-1}$. Frequentist and Bayesian inference now give different
answers (although in practice the answers will be quite similar) and
with very different justifications. This is explained below.

\subsection{Frequentist Inference}

We start with the least squares estimate of $\hat{\beta}$:
\begin{align*}
  \hat{\beta} = (X^T X)^{-1} X^T y
\end{align*}
where
\begin{equation*}
  y =
  \begin{pmatrix}
    y_2 \\ \cdot \\ \cdot \\ \cdot \\ y_n
  \end{pmatrix} ~~~ X =
  \begin{pmatrix}
    1  & y_1\\
    1 & y_2 \\
    \cdot & \cdot \\
    \cdot & \cdot \\
    \cdot & \cdot \\
    1 & y_{n-1} 
  \end{pmatrix} ~~~
  \beta =
  \begin{pmatrix}
    \beta_0 \\ \beta_1 
  \end{pmatrix} ~~~
  \hat{\beta} =
  \begin{pmatrix}
    \hat{\beta}_0 \\ \hat{\beta}_1
  \end{pmatrix}
\end{equation*}
Unlike in the previous case, we cannot compute the distribution of
$\hat{\beta}$ assuming that $X$ is fixed. This is because $X$ involves
depends on $y_1, \dots, y_{n-1}$. Instead, we need to use other
arguments. For example, note that
\begin{align*}
  \hat{\beta}_1 = \frac{\sum_{t=2}^n (y_t - \bar{y})(x_t -
  \bar{x})}{\sum_{t=2}^n (x_t - \bar{x})^2} = \frac{\sum_{t=2}^n (y_t
  - \bar{y}_{2:n})(y_{t-1} - 
  \bar{y}_{1:(n-1)})}{\sum_{t=2}^n (y_{t-1} - \bar{y}_{1:(n-1)})^2}  
\end{align*}
where $\bar{y} = \bar{y}_{2:n} := (y_2 + \dots + y_n)/(n-1)$ and
$\bar{x} = \bar{y}_{1:(n-1)} := (y_1 +
\dots + y_{n-1})/(n-1)$. The distribution of $\hat{\beta}_1$ is
actually quite complicated. Asymptotic arguments can be used to derive
a limiting normal distribution. For the full argument, see, for
example, \citet[Subsection B.3]{ShumwayStoffer} which uses a dependent
Central Limit Theorem.

\subsection{Bayesian Inference}
For Bayesian inference, we need to write the prior and likelihood. The
prior will obviously be the same as in usual linear regression:
\begin{align*}
  \beta_0, \beta_1, \log \sigma \overset{\text{i.i.d}}{\sim}
  \text{unif}(-\infty, \infty). 
\end{align*}
For the likelihood, we write:
\begin{align*}
  \text{Likelihood} &= f_{y_1, \dots, y_n \mid \beta, \sigma}(y_1,
                      \dots, y_n) \\
  &= f_{y_1 \mid \beta, \sigma}(y_1) \prod_{t=2}^{n} f_{y_t \mid y_1,
    \dots, y_{t-1}, \beta, \sigma}(y_t) \\
  &= f_{y_1 \mid \beta, \sigma}(y_1) \prod_{t=2}^n f_{\beta_0 +
    \beta_1 y_{t-1} + \epsilon_t \mid y_1, \dots, y_{t-1}, \beta,
    \sigma} (y_t) \\
  &= f_{y_1 \mid \beta, \sigma}(y_1) \prod_{t=2}^n f_{\epsilon_t \mid
    y_1, \dots, y_{t-1}, \beta, \sigma} (y_t - \beta_0 - \beta_1
    y_{t-1}). 
\end{align*}
We can now assume that $\epsilon_t$ is independent of $y_1, \dots,
y_{t-1}$ for each $t = 2, \dots, n$. This gives
\begin{align*}
  \text{Likelihood} &= f_{y_1 \mid \beta, \sigma}(y_1) \prod_{t=2}^n
                      f_{\epsilon_t}(y_t - \beta_0 - \beta_1 y_{t-1}) \\
  &= f_{y_1 \mid \beta, \sigma}(y_1) \prod_{t=2}^n
    \frac{1}{\sigma\sqrt{2 \pi}} \exp \left(-\frac{(y_t - \beta_0 -
    \beta_1 y_{t-1})^2}{2 \sigma^2} \right). 
\end{align*}
For $f_{y_1 \mid \beta, \sigma}(y_1)$, the simplest thing is to assume
that it does not depend on the parameters $\beta, \sigma$ so that this
term can be ignored. This assumption leads to
\begin{align}
  \text{Likelihood} &= f_{y_1 \mid \beta, \sigma}(y_1) \prod_{t=2}^n
    \frac{1}{\sigma\sqrt{2 \pi}} \exp \left(-\frac{(y_t - \beta_0 -
    \beta_1 y_{t-1})^2}{2 \sigma^2} \right) \nonumber \\ &\propto \sigma^{-(n-1)}
                      \exp \left(-\frac{1}{2 \sigma^2} \|y - X
                      \beta\|^2 \right). \label{like_ar}
\end{align}
This is exactly the same likelihood as in usual linear regression (see
\eqref{liklin}) except that $n$ in \eqref{liklin} is now replaced by
$n-1$ which is the number of observations in the regression model
\eqref{ar1}. Because the likelihood is of the same form as in
\eqref{liklin}, the posterior of $\beta$ follows the same calculation
as in usual linear regression and we obtain:
\begin{align*}
  \beta_0, \beta \mid \text{data} \sim t_{2} \left(\hat{\beta},
  \frac{S(\hat{\beta})}{n-3} (X^T X)^{-1}, n-3 \right) 
\end{align*}
This is the same as \eqref{post3} with $n-1$ replacing $n$ and $m =
2$.

Therefore, Bayesian inference gives the same posterior
$t$-distribution in AutoRegression as in usual linear
regression. Unlike frequentist AutoRegression, there is no need for
asymptotics assuming $n \rightarrow \infty$. This example is useful
for highlighting the differences between Bayesian and frequentist
inference. 




\begin{thebibliography}{99}

%\bibitem[MacKay and Peto(1995)]{MacKay1995HierarchicalDirichlet}
%D.~J.~C. MacKay and L.~C. Peto,
%\newblock ``A hierarchical Dirichlet language model,''
%\newblock \emph{Natural Language Engineering}, vol.~1,
%no.~3, pp.~289--308, 1995.

%\bibitem[Rubin(1981)]{Rubin1981BayesianBootstrap}
%D.~B. Rubin,
%\newblock ``The Bayesian bootstrap,''
%\newblock \emph{The Annals of Statistics}, vol.~9,
%no.~1, pp.~130--134, 1981.
  
%\bibitem[Fisher(1934)]{Fisher1934}
%R.~A. Fisher,
%\newblock ``Probability, likelihood and quantity of information in the logic of
%uncertain inference,''
%\newblock \emph{Proceedings of the Royal Society of London. Series~A,
%Containing Papers of a Mathematical and Physical Character}, vol.~146,
%no.~856, pp.~1--8, 1934.

%\bibitem[Efron(2010)]{Efron}
%Efron, B. (2010)
%\textit{Large-scale inference: empirical Bayes methods for estimation,
%testing, and prediction}. Cambridge University Press, 2010.

\bibitem[Shumway and Stoffer(2017)]{ShumwayStoffer}
Shumway, R. H. and Stoffer, D. S. (2017)
\textit{Time Series Analysis and Its Applications: With R
  Examples}. Springer, 4th edition. 
  
%     \bibitem[Jaynes(2003)]{Jaynes} Jaynes, E. T, (2003)
%  \textit{Probability theory: the logic of science}. Cambridge
%  University Press, 2003.
 
%  \bibitem[Sinai(1992)]{Sinai} Sinai, Y. G, (1992)
%  \textit{Probability theory: an introductory course}. Springer
%  Verlag, 1992.  

%\bibitem[{deGroot(1986)}]{DeGrootBlackwell}DeGroot, M. H. (1986)
%       A conversation with David Blackwell.
%\newblock \emph{Statistical Science}, \textbf{1}, 40--53.

%         \bibitem[Mosteller(1987)]{Mosteller} Mosteller, F, (1987)
%  \textit{Fifty challenging problems in probability with
%    solutions}. Courier Corporation, 1987.

%    \bibitem[MacKay(2003)]{MacKay} MacKay, D. J. C., (2003)
%  \textit{Information theory, inference, and learning
%  algorithms}. Cambridge University Press, 2003.

%\bibitem[{Lindley(2013)}]{Lindley}Lindley, D. V. (2013)
%   \textit{Understanding Uncertainty}. John Wiley and sons, 2013.


\end{thebibliography}







\end{document}

%%% Local Variables:
%%% mode: latex
%%% TeX-master: t
%%% End:
